\documentclass[11pt]{amsproc}

\usepackage[a4paper,margin=2.6cm]{geometry}
\usepackage{amsmath,amssymb,amsthm}
\usepackage[spanish]{babel}   
\usepackage[T1]{fontenc}
\usepackage[utf8]{inputenc}

\newtheorem*{theorem}{Teorema}
\newcommand{\Q}{\mathbb{Q}}
\newcommand{\C}{\mathbb{C}}
\newcommand{\Z}{\mathbb{Z}}

\begin{document}

\section*{Problema}

Una red consta de un número finito de puntos, unidos por flechas. Entre dos
puntos puede haber varias flechas, y también puede haber flechas que comiencen
y terminen en un mismo punto.
Un viaje consiste en seguir sucesivamente algunas flechas. Sea $N_k$ el
número de viajes de $k$ pasos cuyo punto final coincide con el inicial.
Demuestre que, si $p$ es primo, entonces
\[
N_p\equiv N_1\bmod p.
\]

\section*{Primera solución}

Consideremos el conjunto $X$ de los viajes cerrados de $p$ pasos. El grupo
cíclico $C_p$ actúa sobre $X$ rotando los pasos del viaje: si
$(e_1,e_2,\dots,e_p)$ 
es un viaje cerrado, lo enviamos a
$(e_2,e_3,\dots,e_p,e_1)$.
Como $p$ es primo, toda órbita de esta acción tiene 1 o $p$ elementos.
Por lo tanto,
$|X|\equiv |X^{C_p}|\bmod p$, 
donde $X^{C_p}$ denota el conjunto de viajes fijos por la rotación.

Un viaje es fijo si y solo si sus $p$ pasos son una misma flecha repetida.
Para que esto sea posible, esa flecha debe comenzar y terminar en el mismo
punto. En consecuencia, los viajes fijos están en correspondencia con los
viajes cerrados de un solo paso. Así,
$|X^{C_p}|=N_1$. 
Como $|X|=N_p$, obtenemos
$N_p\equiv N_1\bmod p$.

\section*{Otra solución, más algebraica}

Sea $A=(a_{ij})$ la matriz de adyacencia de la red, donde $a_{ij}$ es el
número de flechas que van del punto $i$ al punto $j$. Es bien conocido, y
se comprueba por inducción, que la entrada $(i,j)$ de $A^k$ cuenta los
viajes de $k$ pasos que comienzan en $i$ y terminan en $j$. Por lo tanto,
\[
N_k=\operatorname{tr}(A^k).
\]

Usaremos la siguiente versión matricial del pequeño teorema de Fermat.

\begin{theorem}
Sea $A$ una matriz cuadrada con coeficientes enteros y sea $p$ un número
primo. Entonces
\[
\operatorname{tr}(A^p)\equiv
\bigl(\operatorname{tr}A\bigr)^p\bmod p.
\]
\end{theorem}

\begin{proof}
Sean $\lambda_1,\dots,\lambda_n\in\mathbb C$ los autovalores de $A$,
contados con multiplicidad. Como $A$ tiene coeficientes enteros, todos ellos
son enteros algebraicos.
Al desarrollar
\[
(\lambda_1+\cdots+\lambda_n)^p,
\]
todos los términos mixtos tienen coeficiente divisible por $p$. Por tanto
existe un entero algebraico $\alpha$ tal que
\[
(\lambda_1+\cdots+\lambda_n)^p
\lambda_1^p+\cdots+\lambda_n^p+p\alpha.
\]
Como 
\[
\lambda_1+\cdots+\lambda_n=\operatorname{tr}A,\quad 
\lambda_1^p+\cdots+\lambda_n^p=\operatorname{tr}(A^p),
\]
tenemos 
\[
p\alpha=
\bigl(\operatorname{tr}A\bigr)^p-\operatorname{tr}(A^p)\in\mathbb Z.
\]
En particular, $\alpha\in\mathbb Q$. Como $\alpha$ es además un entero
algebraico, necesariamente $\alpha\in\mathbb Z$. Luego 
\[
\operatorname{tr}(A^p)\equiv
\bigl(\operatorname{tr}A\bigr)^p\bmod p.\qedhere 
\]
\end{proof}

Volvamos a la solución del problema. 
Aplicando el teorema a la matriz de adyacencia, 
\[
N_p=\operatorname{tr}(A^p)
\equiv \bigl(\operatorname{tr}A\bigr)^p
=N_1^p\bmod p.
\]
Por el pequeño teorema de Fermat,
$N_1^p\equiv N_1\bmod p$, y entonces 
$N_p\equiv N_1\bmod p$. 

\end{document}
